\chapter{Formalizing the semantics in Coq}
\label{chap:formalization}

\section{Representing syntax, heaps, and judgments}
% Explain the mutually inductive syntax and evaluation relations, finite-map
% heaps, and list representations of argument stacks and black-hole sets.
% Describe which syntax restrictions the relations enforce rather than the types.

\section{Binding, substitution, and local definitions}
% Explain variable-pattern substitution and allocation for where bindings.
% Compare the implemented substitution and renaming operations with the source
% paper's binding conventions. State the qualifications of the current model.

\section{Choosing fresh names}
\label{sec:freshness}
% Use a short where example to show why the two semantics must choose compatible
% names when correspondence requires identical heaps and results.

\subsection{Canonical names and the counter-based approach}
% Explain the alpha-equivalence obstacle in the first correspondence attempt.
% Describe how counter threading aligned names and what it added to judgments.
% Keep the history focused on the proof obligation each approach addressed.

\subsection{A shared fresh-name generator}
% Present FreshVarGen and its laws, then explain how both semantics use it.
% Distinguish agreement on generated names from the source paper's full
% freshness condition: the current calls avoid L or update_locs St, not every
% name occurring in the heap, arguments, and matching. Note that main supplies
% no concrete generator instance.

\section{Determinism}
% State determinism for expression evaluation, matching evaluation, and machine
% transitions with a fixed generator. Explain the role of arity and fresh-name
% choices without reproducing the tactic scripts.
